% 底本：EGAIV.1.pdf，PDF第168--170页，原分册印刷第169--171页。
\subsection{Cartier等式}
下面的结果把MacLane关于导子的一个定理用微分表述出来。

\egas{21.7.1}{定理（Cartier）}
设 $K$ 为域，$L/K$ 为有限型扩张。
则 $\Omega^1_{L/K}$ 与 $\Upsilon_{L/K}$ 都是有限维 $L$-向量空间，并且
\begin{equation}
 \rk_L\Omega^1_{L/K}-\rk_L\Upsilon_{L/K}=\trdeg_KL.
 \tag{21.7.1.1}
\end{equation}

\eproof
若 $K\subseteq L'\subseteq L$，则由（20.6.17.1）
（取 $\Lambda=P$ 为 $K$ 的素域、$A=K$、$B=L'$、$C=L$），有正合列
\begin{equation*}
\begin{gathered}
 0\longrightarrow\Upsilon_{L'/K}\otimes_{L'}L
 \longrightarrow\Upsilon_{L/K}\longrightarrow\Upsilon_{L/L'}\\
 \longrightarrow\Omega^1_{L'/K}\otimes_{L'}L
 \longrightarrow\Omega^1_{L/K}\longrightarrow\Omega^1_{L/L'}
 \longrightarrow0.
\end{gathered}
\end{equation*}
\cmt{L12-A01}{原文引（20.6.15.1）。作者在《勘误与增补（清单三）》A部正式改为（20.6.17.1）；正文从其修订。
这里的六项正合列，而不是单纯的复形定义，是证明所需要的依据。}
当其中各项有限维时，由（$0_{\mathrm{III}}$，11.10.2）得到
\begin{align*}
 \rk_L\Omega^1_{L/K}-\rk_L\Upsilon_{L/K}
 &=\bigl(\rk_L\Omega^1_{L/L'}-\rk_L\Upsilon_{L/L'}\bigr)\\
 &\quad+\bigl(\rk_{L'}\Omega^1_{L'/K}-\rk_{L'}\Upsilon_{L'/K}\bigr).
\end{align*}
又有 $\trdeg_KL=\trdeg_KL'+\trdeg_{L'}L$。
所以沿有限个生成元递推，只须处理 $L=K(x)$ 的情形。
各项的有限维性也可由同一正合列同时递推得到：若左右两段扩张的两个模有限维，
则中间扩张的两个模均由有限维空间的子商和扩张构成，因而有限维。
\cmt{L12-G06}{原文将有限维性隐含在“对生成元数递推”的论证中。
正文明确把它与等式一同递推，随后才使用有限正合列的Euler特征；
这避免在尚未确定有限维时书写可能无意义的“无穷减无穷”。}
区分三种情形。

\noindent a) $x$ 在 $K$ 上超越。
$L/K$ 可分，故 $\Upsilon_{L/K}=0$（20.6.3）；
由（20.5.9）和（20.4.13），$\Omega^1_{L/K}$ 的维数为1，故等式成立。

\noindent b) $L/K$ 是有限可分代数扩张。
此时仍有 $\Upsilon_{L/K}=0$（20.6.3），且 $\Omega^1_{L/K}=0$（20.6.20），
所以等式仍成立。

\noindent c) $L/K$ 是不可分代数扩张。
利用开头的递推，可以归结为 $L^p\subseteq K$ 的情形。
由（21.4.8），有
$\rk_L\Omega^1_{L/K}=\rk_L\Upsilon_{L/K}$，故仍得（21.7.1.1）。
\cmt{L12-G07}{对简单不可分扩张，可先令 $x^{p^e}$ 在 $K$ 上可分，
再把 $K(x)/K(x^{p^e})$ 分成若干指数一的纯不可分扩张；各级均由一个元素生成。
指数一简单扩张的 $p$-基至多含一个元素，因此（21.4.5）、（21.4.8）同时给出所需有限维性。
这里补出原文“可归结为”的具体过程，不增加定理的适用假设。}
\finishproof

\egas{21.7.2}{推论}
设 $K$ 为域，$L/K$ 为有限型扩张，$k$ 为 $K$ 的子域。
则 $\Omega^1_{L/K}$、$\Upsilon_{L/K/k}$ 均为有限维 $L$-向量空间，并且
\begin{equation}
 \rk_L\Omega^1_{L/K}-\rk_L\Upsilon_{L/K/k}
 =\trdeg_KL+d(L/K,k).
 \tag{21.7.2.1}
\end{equation}
因此
\begin{equation}
 \rk_L\Omega^1_{L/K}-\rk_L\Upsilon_{L/K/k}\ge\trdeg_KL.
 \tag{21.7.2.2}
\end{equation}
等号成立当且仅当 $k$ 是扩张 $L/K$ 的可容许域。

事实上，$\Upsilon_{L/K}\to\Upsilon_{L/K/k}$ 满射；按（21.6.1.2），
\[
 \rk_L\Upsilon_{L/K}-\rk_L\Upsilon_{L/K/k}=d(L/K,k).
\]
于是结论来自（21.7.1）和（21.6.2）。

\egas{21.7.3}{推论}
设 $k\subseteq K\subseteq L$ 为三个域，且 $L/k$ 为有限型扩张。
则
\begin{equation}
 \rk_L\Omega^1_{L/k}-\rk_K\Omega^1_{K/k}
 =\trdeg_KL+d(L/K,k).
 \tag{21.7.3.1}
\end{equation}
从而
\begin{equation}
 \rk_L\Omega^1_{L/k}-\rk_K\Omega^1_{K/k}\ge\trdeg_KL,
 \tag{21.7.3.2}
\end{equation}
等号成立当且仅当 $k$ 是 $L/K$ 的可容许域。

由正合列（20.6.1.1）
\[
 0\longrightarrow\Upsilon_{L/K/k}
 \longrightarrow\Omega^1_{K/k}\otimes_KL
 \longrightarrow\Omega^1_{L/k}\longrightarrow\Omega^1_{L/K}
 \longrightarrow0,
\]
（21.7.3.1）左端等于
$\rk_L\Omega^1_{L/K}-\rk_L\Upsilon_{L/K/k}$，再应用（21.7.2）即可。

这个推论的好处是只涉及微分模，而不再显式涉及非完美性模。
下文（21.8.6）还将证明：对任意有限型扩张 $L/K$，
都存在 $K$ 的子域 $k$，使 $[K:k]$ 有限，且 $k$ 是 $L/K$ 的可容许域；
此时（21.7.3.2）两端相等。

\egas{21.7.4}{推论}
设 $K$ 为域 $k$ 的有限型扩张。下列条件等价：
\begin{enumerate}[label=\alph*),leftmargin=2.5em]
\item $K/k$ 为有限可分扩张；
\item $K$ 是形式非分歧的 $k$-代数（19.10.2）；
\item $K$ 是形式平展的 $k$-代数（19.10.2）；
\item $\Omega^1_{K/k}=0$。
\end{enumerate}

\eproof
所有拓扑均取离散拓扑。b)与d)由（20.7.4）等价；
a)蕴涵形式光滑性（19.6.1），结合形式非分歧性即得到c)。
另一方面c)蕴涵b)。所以只须证明d)推出a)。
由（21.7.1.1），$\Omega^1_{K/k}=0$ 迫使 $\trdeg_kK=0$ 且 $\Upsilon_{K/k}=0$。
前者与有限型假设一起表明扩张有限，后者由（20.6.3）表明它可分。
\finishproof

\egas{21.7.5}{注记}
由（20.6.5.2）和（$0_{\mathrm{III}}$，11.10.2），
（21.7.1.1）的左端就是复形 $K_\bullet(C/B/A)$ 的Euler--Poincaré特征，
其中 $C=L$、$B=K$、$A=P$ 为 $K$ 的素域；该复形定义于（20.6.5）。
在本专著关于相交理论与Riemann--Roch定理的章节中，
推广（20.6.26）所讨论的复形 $F_\bullet(C/A)$ 后得到的复形，
其取值于模的类群的广义Euler--Poincaré特征也将发挥重要作用。
\cmt{L12-T04}{原文依次引（20.6.3.2）、（20.6.3）和（20.6.22）；
对应的同调公式、二项复形定义以及 $F_\bullet(C/A)$ 的讨论实际在
（20.6.5.2）、（20.6.5）和（20.6.26）。正文按内容定位更正。
这里的Euler特征是有限维同调的交错和；不要求构成复形的两个项本身有限维。}
